\section{A stronger multiplicity exclusion from projections}

The following two findings give a comparison of actual coordinate-ring
multiplicities, rather than just nonvanishing of an upper bound.
They are derived here in an explicit form; no assertion of historical
priority is intended. They use the polynomial-functor dictionary of
cote~133, scan~028, but require an additional determinantal construction.

\subsection*{Finding M-006: projection containment compares all copies}
Let \(V\) be a complex vector space, let \(U\) be an \(r\)-dimensional
quotient of \(V\), and choose a linear section \(U\hookrightarrow V\).
Write
\[
 F_d(V)=\operatorname{Sym}^d(\operatorname{Sym}^n V),\qquad
 \pi:\operatorname{Sym}^n V^*\longrightarrow
      \operatorname{Sym}^n U^*
\]
for restriction followed, when necessary, by the inclusion of forms
depending only on \(U\).
Let \(X,Y\subseteq\operatorname{Sym}^n V^*\) be reduced closed
\(\mathrm{GL}(V)\)-stable cones. Suppose
\begin{equation}\label{eq:M006hyp}
                         \pi(Y)\subseteq X.
\end{equation}
Then for every \(d\) and every partition \(\lambda\vdash nd\)
with \(\ell(\lambda)\leq r\),
\begin{equation}\label{eq:M006}
 \operatorname{mult}_{S_\lambda V}\mathbb C[Y]_d
 \ \leq\
 \operatorname{mult}_{S_\lambda V}\mathbb C[X]_d.
\end{equation}
Only the stated containment of the projected variety is required;
containment of \(Y\) itself has not been assumed.

\begin{proof}
Choose coordinates in which \(U\) is the span of the first \(r\)
basis vectors. The natural Schur--Weyl decomposition gives
\[
 F_d(V)=\bigoplus_{\lambda\vdash nd} S_\lambda V\otimes M_\lambda,
 \qquad
 M_\lambda=[\lambda]^{S_n\wr S_d}.
\]
The ideals \(I(X)_d\) and \(I(Y)_d\) are
\(\mathrm{GL}(V)\)-submodules; hence, in characteristic zero, they
have the forms
\[
 I(X)_d=\bigoplus_\lambda S_\lambda V\otimes J_{X,\lambda},
 \qquad
 I(Y)_d=\bigoplus_\lambda S_\lambda V\otimes J_{Y,\lambda}
\]
for subspaces of the same \(M_\lambda\).
The natural inclusion \(F_d(U)\hookrightarrow F_d(V)\) identifies
\(S_\lambda U\otimes M_\lambda\) with the corresponding summand
of \(F_d(U)\). If \(\ell(\lambda)\leq r\), it is nonzero.
Any \(f\in S_\lambda U\otimes J_{X,\lambda}\) belongs to \(I(X)_d\)
and depends only on the projection to \(U\). For \(y\in Y\),
\[
                              f(y)=f(\pi(y))=0
\]
by \eqref{eq:M006hyp}. Since \(Y\) is reduced over the algebraically
closed field \(\mathbb C\), this means \(f\in I(Y)_d\).
The \(\mathrm{GL}(V)\)-span of the nonzero
\(S_\lambda U\subseteq S_\lambda V\) is the irreducible
\(S_\lambda V\). Thus
\(S_\lambda V\otimes J_{X,\lambda}\subseteq I(Y)_d\), giving
\(J_{X,\lambda}\subseteq J_{Y,\lambda}\).
Subtracting their dimensions from \(\dim M_\lambda\) proves
\eqref{eq:M006}. This proves the comparison for all copies of the
representation, not just its presence.
\end{proof}

\paragraph{Equivalent highest-weight formulation.}
A polynomial highest-weight vector of weight
\((\lambda_1,\ldots,\lambda_r,0,\ldots,0)\) lies in \(F_d(U)\):
the polynomial weights have nonnegative coordinates, so trailing
zero weights allow no coordinate outside \(U\).
The proof therefore says that every equation of type \(\lambda\)
for \(X\) is also an equation for \(Y\) once all its projected
arguments already lie in \(X\).

\subsection*{Finding M-007: an explicit determinant for every padded form}
Let \(m\geq2\), \(r\geq2\), and put
\[
 p=\left\lfloor\frac{r-1}{2}\right\rfloor,\quad
 q=\left\lceil\frac{r-1}{2}\right\rceil,\qquad
 S(m,r)=\binom{m+p}{p}+\binom{m+q}{q}-1.
\]
For any degree-\(m\) homogeneous form \(f\) in at most \(r\)
variables and any nonzero linear form \(L\) in the same variable
space, the form \(L^{n-m}f\) is the determinant of an \(n\times n\)
matrix of homogeneous linear forms whenever \(n\geq S(m,r)\).
The entries may use the coefficients of \(f\) as constants.
The underlying determinant identity is an identity over
\(\mathbb Z\) in universal coefficient variables.

\begin{proof}
Take \(L\) as one coordinate and split the other \(r-1\) coordinates
into \(x=(x_1,\ldots,x_p)\) and \(y=(y_1,\ldots,y_q)\). Write
\begin{equation}\label{eq:M007expansion}
 f=\sum_{|\alpha|+|\beta|\leq m}
       c_{\alpha,\beta}x^\alpha y^\beta
                      L^{m-|\alpha|-|\beta|}.
\end{equation}
For \(p=0\) the only \(x\)-multi-index is the empty one.

Construct a rooted tree whose vertices are all multi-indices
\(\alpha\in\mathbb Z_{\geq0}^{p}\) with \(|\alpha|\leq m\).
Its root is \(0\). For every nonzero \(\alpha\), let \(i\) be
the largest index for which \(\alpha_i>0\), and draw the edge
\(\alpha-e_i\longrightarrow\alpha\), with label \(-x_i\).
There is a unique path from the root to \(\alpha\); its product
of labels is \((-1)^{|\alpha|}x^\alpha\).
Make a second tree on the analogous \(y\)-multi-indices
\(\beta\), but reverse all its edges, using label \(-y_j\).
Its unique path from \(\beta\) to the root has product
\((-1)^{|\beta|}y^\beta\).
Identify the two roots into one central vertex, and keep all
other vertices distinct.

For each pair with \(|\alpha|+|\beta|\leq m\), insert a directed
edge from the prefix vertex \(\alpha\) to the suffix vertex
\(\beta\), labelled \(c_{\alpha,\beta}L\).
The pair \((0,0)\) gives the central loop
\(c_{0,0}L\). Put a loop labelled \(L\) at every noncentral
vertex. Let \(M\) be the matrix whose entry \(M_{uv}\) is the
sum of the labels of all edges \(u\to v\).
All its entries are homogeneous linear forms, and its size is
the number of vertices,
\[
                    S=\binom{m+p}{p}+\binom{m+q}{q}-1.
\]

We now evaluate the determinant by its permutation expansion.
Aside from loops, every directed cycle must pass through the
central vertex. Indeed, the prefix tree only points outwards,
the suffix tree only points inwards, and coefficient edges
only point from a prefix vertex to a suffix vertex.
A cycle through the central vertex uses exactly one coefficient
edge, from some \(\alpha\) to some \(\beta\).
Consequently a cycle cover either uses the central loop, or
chooses exactly one such nontrivial cycle and uses loops at
all unused vertices.

Put \(k=|\alpha|+|\beta|\). The chosen cycle has \(k+1\) edges:
\(k\) variable edges and one coefficient edge.
Its permutation sign is \((-1)^k\); its variable-edge labels
have the same sign. The signs cancel. The coefficient edge
and the \(S-k-1\) unused loops together contribute \(L^{S-k}\).
This cycle cover therefore contributes exactly
\[
                   c_{\alpha,\beta}x^\alpha y^\beta L^{S-k}.
\]
The central-loop case gives the same expression for \(k=0\).
Summing all cycle covers and using \eqref{eq:M007expansion} gives
\[
                                \det M=L^{S-m}f.
\]
For \(r\geq2\), \(S\geq m+1\), so the exponent is nonnegative.
If \(n>S\), adjoining a diagonal block \(L I_{n-S}\) gives the
required matrix of size \(n\).
All cancellations and coefficient signs in this argument are
integral polynomial identities, so the construction with
indeterminate \(c_{\alpha,\beta}\) is defined over \(\mathbb Z\).
\end{proof}

\paragraph{Boundary cases and complexity meaning.}
For \(r=1\), \(f\) is a scalar multiple of \(L^m\), and a
diagonal determinant suffices. A degree-one form is similarly
immediate. If \(L=0\) and \(n>m\), the padded form is zero and
has a zero determinantal representation.
The result is a size bound with arbitrary scalar coefficients;
it does not assert that the coefficients of a projected
permanent can be computed in polynomial time.

\subsection*{Finding M-008: row bounds for actual GCT multiplicity obstructions}
Let \(m\geq2\), \(n>m\), \(V=\mathbb C^{n^2}\), and use
\(q^m_\lambda(d[n])\) and \(K_n(\lambda)\) for the actual
padded-permanent and determinant-closure multiplicities.
If \(n\geq S(m,r)\) and \(r\leq n^2\), then
\begin{equation}\label{eq:M008}
        q^m_\lambda(d[n])\leq K_n(\lambda)
   \qquad\text{for every \(d\) and every \(\ell(\lambda)\leq r\)}.
\end{equation}
In particular:
\begin{enumerate}
\item If \(n>3m^4\), no actual multiplicity obstruction has
      at most nine rows.
\item If \(n\geq m^{25}\), no actual multiplicity obstruction
      has at most fifty-one rows.
\item More generally, for every fixed integer \(c\geq6\),
      \(n\geq m^c\) excludes all outputs with at most \(2c+1\)
      rows.
\end{enumerate}
These are statements about the closure multiplicity \(K_n\),
so they also exclude the stronger sufficient test
\(q^m_\lambda>g_{n,d}(\lambda)\) in these ranges.

\begin{proof}
An \(r\)-variable restriction of any translate of the padded
permanent has the form \(L^{n-m}f\), where \(f\) has degree \(m\)
in those variables. Finding M-007 makes it a determinant when
\(L\ne0\); the zero case was handled above.
Every linear substitution into \(\det_n\) belongs to
\(\Omega_n\), since \(\mathrm{GL}(V)\) is dense in
\(\operatorname{End}(V)\). Thus the projection of the entire
padded-permanent orbit lies in \(\Omega_n\).
The preimage of this closed set under the restriction morphism
is closed, so the same containment holds for the orbit closure.
Apply M-006 with \(Y=Z_{n,m}\) and \(X=\Omega_n\).

For the numerical bounds, when \(r=2c+1\) one has
\(p=q=c\) and \(S(m,r)=2\binom{m+c}{c}-1\).
For fixed \(c\), the real function
\[
 \frac{2\binom{m+c}{c}}{m^c}
       =\frac{2}{c!}\prod_{j=1}^c\left(1+\frac jm\right)
\]
is decreasing for \(m>0\).
For \(m\geq2\) it is at most
\[
                  \frac{(c+1)(c+2)}{2^c}.
\]
For \(c\geq6\), this is at most one: it is \(56/64\)
at \(c=6\), and the ratio at the next integer is
\((c+3)/(2(c+1))<1\).
Therefore \(S(m,2c+1)<m^c\), proving the third assertion
and, with \(c=25\), the second.
For \(c=4\), the same bound is \(30/16=15/8<3\);
thus \(S(m,9)<3m^4\), proving the first.
All these row counts fit inside \(n^2\) in the indicated ranges.
\end{proof}

\paragraph{Exact scope of the new gate.}
At the Ikenmeyer--Panova threshold \(n>3m^4\), the remaining
actual-multiplicity candidates have at least ten rows; at the
B\"urgisser--Ikenmeyer--Panova threshold \(n\geq m^{25}\), they
have at least fifty-two. They must also have positive
permanent-side multiplicity and satisfy the usual padded-shape
conditions. No assertion that these remaining candidates are
obstructions has been proved.
Every ambient constituent of \(\operatorname{Sym}^d\operatorname{Sym}^nV\)
has at most \(d\) rows: it is a constituent of \(h_n^d\) under the
character map, and Pieri's rule adds at most one new row at each step
(Macdonald, I, (5.16)). Thus the respective remaining actual-obstruction
degrees also satisfy \(d\geq10\) and \(d\geq52\).
This is independent of the stronger inequality \(d>n/m\), which is
known to be necessary for the sufficient test \(q>g\), not for every
possible \(q>K_n\).
The four-row ambient loss in O-008 is consistent with M-008:
retaining full ambient multiplicity and retaining at least the
permanent multiplicity are different assertions.

\paragraph{Grothendieck dictionary.}
The restriction map and the Schur decomposition in M-006 are
polynomial-functor operations on locally free sheaves, as in
cote~133, scan~028, and Macdonald, I, Appendix~A,
(5.3)--(5.6), pp.~155--156.
The matrix in M-007 defines a morphism over
\(\mathbb Z[c_{\alpha,\beta}]\) from an affine space of
coefficient data to the affine scheme of linear matrix pencils.
Its determinant is the displayed universal section of
\(\mathcal O(n)\). The projection containment is a containment
of reduced characteristic-zero orbit-closure schemes, and its
ideal inclusion on isotypic summands is the proof of M-008.
SGA~3, I, 4.7.2 supplies the equivariant-module language and
SGA~6, VI, 3.3 the lambda operations on their \(K_0\).
No semisimple splitting over all of \(\mathbb Z\) is inferred
from the characteristic-zero multiplicity comparison.
